\chapter{Ethical Considerations}
Computational drug discovery is the bridge between machine learning and human health, where errors or misuses can have real-world consequences.
DTI predictions would be the driving force for drug development pipelines, so a false positive would waste research resources, but a false negative can overlook a promising drug candidate.


\section{Data Licensing}
\subsection{BindingDB}
This work uses data from the BindingDB database, downloaded in November 2025.
BindingDB is a publicly accessible where the data is curated directly by BindingDB staff and released under the Creative Commons Attribution 3.0 (CC BY 4.0) licence.
The licence permits use, redistribution, and adaptation of the data for non-commercial research purposes, provided appropriate attribution is given.
This work cites BindingDB accordingly and uses the data solely for academic research.
\subsection{ProtBERT}
For protein sequence encoding, this project uses ProtBERT, a pre-trained protein language model developed by Rostlab and distributed via Hugging Face.
The ProtTrans model family, including ProtBERT, is released under the Academic Free Licence v3.0 (AFL-3.0), which permits free use, modification, and redistribution for academic purposes.
The model weights were used as provided, without modification to the pre-trained parameters, and the original work is cited in this dissertation.

The use of both dataset and pre-trained models is compliant with their respective licensing terms, and no restricted data was used without permission.
\section{Model Misuse}
This tool is purely for academic purposes and not clinical decision-making, therefore, a certain tolerance level needs to be considered for model, biological, and chemical accuracy.
\subsection{Easy Accessibility}
\label{subsec:easy-accessibility}
By wrapping our model in an intuitive web interface, it is making virtual screening accessible to more people who might not have pharmacological knowledge.
Users might not fully understand the limitations and real meaning of the predictions, so there's an ethical obligation to communicate uncertainty clearly.
The barrier for misuse has been lowered, through greater accessibility.
\subsection{Dual Use}
Another issue with broadening the type of users who have access to the product is they may twist the model the wrong way
The same system that can identify therapeutic drug-target interactions can theoretically be used to identify toxic interactions.
This is only a university project, however, it is important to consider all pathways if a similar product were to go onto market.
A way to limit this as it is a concern for chemical biological security, is to only show bindings that are successful.
\subsection{Over-Reliance}
If a user treats this DTI tool as a definitive rather than a prioritisation tool, it could lead to a significant waste in resources or unsafe compound progressing.
\section{Biological Data Bias}
The BindingDB dataset tends to over-represent certain protein families (Kinase, GPCRs) and underrepresent others, so our model will naturally perform better on more popular targets.
There is also a publication bias, where more positive interactions are going to be published, hence the imbalanced dataset performed better than the more balanced dataset.
The judgement of the binding thresholds when binarizing the data and also producing the results has genuine impacts further along the pipeline, making this work unique compared to others.
